\chapter{Texte~: du sac de mots aux plongements}\label{ch:ml:text-from-bag-of-words-to-embeddings}

Un titre annonce qu'une entreprise a dépassé les estimations, et son action chute. Un dictionnaire classe le titre comme positif~; un modèle
entraîné sur les rendements apprend que «~dépasse les estimations mais abaisse ses prévisions~» est une mauvaise nouvelle~; et aucun des deux ne sait laquelle des trois sociétés
appelées Apex le titre concerne. Le livre~7 (chapitre~12) a présenté le texte comme donnée et son premier outil, le dictionnaire~;
le livre~8 (chapitre~17) a fait commerce de dépêches lisibles par machine. Ce chapitre construit l'outillage qui relie les mots à la prévision~:
comment un document devient un vecteur, comment ce vecteur est évalué par rapport aux rendements, ce que les thèmes et les plongements peuvent et ne peuvent pas
voir, et comment une mention devient une valeur. Sur un corpus synthétique à vérité connue, un modèle supervisé par les rendements atteint
un coefficient d'information de 0.27 là où un dictionnaire atteint 0.22, des plongements appris sans rendements n'atteignent que 0.09 parce
qu'ils placent «~dépasse~» à côté de «~manqué~», et un lieur d'entités qui utilise les tickers du jour associe une mention sur huit à la
mauvaise société, ou à aucune.

\section{Des documents aux vecteurs}

\begin{definition}[Tokenisation, sac de mots, n-gramme, tf-idf]\label{def:ml:text-from-bag-of-words-to-embeddings:bow}
La \emph{tokenisation}\index{tokenisation} découpe un texte en unités (mots, fragments de mots, ponctuation), en général après
normalisation de la casse. Un \emph{sac de mots}\index{sac de mots} représente un document par les effectifs de ses jetons, en oubliant
leur ordre. Un \emph{n-gramme}\index{n-gramme} est une suite de $n$ jetons consécutifs~; les sacs de n-grammes conservent un peu d'ordre.
La \emph{fréquence du terme--fréquence inverse de document}\index{fréquence du terme--fréquence inverse de document} (tf-idf) pondère l'effectif
du terme $w$ dans le document $d$ par $\log(N/n_w)$, où $n_w$ parmi les $N$ documents contiennent $w$, de sorte que les mots
présents partout comptent peu~; les lignes sont ensuite normalisées à norme unitaire.
\end{definition}

Le corpus du chapitre compte 20\,000 titres sur six ans (1\,512 jours de séance) portant sur environ 32 sociétés~: 28 au départ,
dont trois appelées Apex (une société de logiciels, un mineur et une entreprise de santé), six changements de ticker, trois changements de nom et quatre
radiations dont les tickers sont réattribués à de nouvelles cotations 30 jours plus tard. Chaque titre nomme une société (par son nom complet, son
nom abrégé ou son ticker), énonce l'un des 19 types d'événements formulés de plusieurs manières, et mentionne parfois un mot
de secteur~; chacun porte un rendement de réaction égal à l'effet injecté de l'événement, auquel s'ajoute un bruit d'écart-type 3~\%. Les premiers
60~\% des jours (11\,923 titres) servent à l'entraînement~; les derniers 40~\% (8\,077) servent au test.

Un titre tel que «~Apex beat estimates but cut guidance amid copper demand~» traverse trois étapes avant qu'un modèle ne le voie.
La mention de la société est reconnue et remplacée par un marqueur, afin qu'aucun modèle n'apprenne le nom d'une société comme
signal~; le texte est mis en minuscules et découpé en jetons
\texttt{<co> beat estimates but cut guidance amid copper demand}~; les bigrammes ajoutent \texttt{<co> beat}, \texttt{beat estimates}, \texttt{estimates but}, \texttt{but cut},
\texttt{cut guidance} et trois autres. Les 123 mots des titres d'entraînement (ceux qui figurent dans au moins trois titres) deviennent 473
variables avec les bigrammes~; dans les dépêches réelles, le vocabulaire se chiffre en dizaines de milliers, et les bigrammes le multiplient encore.

\section{Sentiment à partir des dictionnaires et des rendements}

Un dictionnaire (livre~7, chapitre~12) est un instrument fixe~: les listes construites à la main dans ce chapitre comptent 11 mots positifs et 16 mots
négatifs. Il lit correctement un résultat («~beat~» obtient en moyenne $+0.20$, pour un effet injecté de $+1~\%$~; une baisse de prévisions
obtient $-0.18$ pour $-1.5~\%$) et se trompe dès qu'il s'agit de plus d'un mot~: «~beat estimates but cut guidance~» obtient
exactement zéro, un coup dans chaque sens, alors que l'effet injecté est de $-1.2~\%$~; «~did not cut guidance~» obtient un score négatif pour un
effet injecté de $+0.3~\%$~; «~wins industry award~» n'obtient rien de positif~; «~announces share buyback~» n'obtient aucun score pour
$+0.7~\%$. Loughran et McDonald (2011) ont montré que les listes de mots générales classent mal les mots financiers~; toute liste rate les
mots que ses auteurs n'ont pas envisagés et ne sait pas pondérer ceux qu'elle contient.

L'alternative consiste à laisser les rendements choisir les mots. Le crible supervisé par les rendements du \cref{lst:ml:text:supervised} conserve les
mots dont la moyenne des rendements des titres qui les contiennent diffère de celle des autres de plus de trois erreurs types, et pondère chacun par cet
écart, dans l'esprit de Ke, Kelly et Xiu (2019), dont le modèle combinant criblage et thèmes construit un score de sentiment adapté
à la prévision des rendements. Sur la moitié d'entraînement, il retient 41 mots, de «~trims~» ($-1.43~\%$) et «~halves~» jusqu'à «~takeover~»
($+1.79~\%$)~; il retient aussi «~to~» et «~be~», qui ne portent de l'information que par la façon dont les titres sur les OPA du corpus
sont formulés, le genre de mot qu'une liste d'exclusion aurait retiré et qu'un crible sur les rendements garde à bon droit.
Le \Cref{tab:ml:text:ic} compare les scores sur la moitié de test~; la \cref{fig:ml:text:learning} montre comment chacun dépend du
nombre de titres étiquetés.

\begin{table}[htbp]
\centering\small
\begin{tabular}{@{}>{\raggedright\arraybackslash}p{11cm}>{\centering\arraybackslash}p{4.6cm}@{}}
\toprule
score & coefficient d'information de rang avec le rendement de réaction\\
\midrule
dictionnaire (11 mots positifs, 16 mots négatifs) & 0.224\\
tf-idf en mots, \omterm{def:ml:linear-and-regularised-baselines:logistic}{régression logistique} & 0.273\\
tf-idf en mots et bigrammes, \omterm{def:ml:linear-and-regularised-baselines:logistic}{régression logistique} & 0.257\\
liste de mots supervisée par les rendements & 0.238\\
\omterm{def:ml:text-from-bag-of-words-to-embeddings:embedding}{plongements de mots} (décomposition en valeurs singulières de la cooccurrence), régression ridge & 0.094\\
règles d'événements, rendement moyen d'entraînement par événement & 0.292\\
\midrule
effet injecté & 0.295\\
\bottomrule
\end{tabular}
\caption{Coefficients d'information de rang sur les 8\,077 titres de test, modèles entraînés sur les 11\,923 titres antérieurs. Données~:
\texttt{ml\_text.scores}.}\label{tab:ml:text:ic}
\end{table}

\begin{omfigure}
\begin{tikzpicture}
\begin{semilogxaxis}[width=11cm,height=6.2cm,xlabel={titres d'entraînement étiquetés},ylabel={coefficient d'information de rang sur la moitié de test},
  tick label style={font=\small},label style={font=\small},xmin=250,xmax=12000,ymin=-0.05,ymax=0.32,
  xtick={300,1000,3000,10000},xticklabels={300,1\,000,3\,000,10\,000},grid=major,grid style={gray!25},legend columns=3,
  legend style={font=\footnotesize,at={(0.5,-0.25)},anchor=north,draw=none,fill=none,
  /tikz/every even column/.append style={column sep=6pt}},table/col sep=comma]
\addplot[black!60,thick,dashed] table[x=docs,y=truth]{figdata/ml/13-text-from-bag-of-words-to-embeddings/learning.csv};
\addplot[omThm,very thick,mark=*] table[x=docs,y=dictionary]{figdata/ml/13-text-from-bag-of-words-to-embeddings/learning.csv};
\addplot[omProp,very thick,mark=square*] table[x=docs,y=tfidf1]{figdata/ml/13-text-from-bag-of-words-to-embeddings/learning.csv};
\addplot[omProp!60,thick,mark=square] table[x=docs,y=tfidf2]{figdata/ml/13-text-from-bag-of-words-to-embeddings/learning.csv};
\addplot[omDef,very thick,mark=triangle*] table[x=docs,y=supervised]{figdata/ml/13-text-from-bag-of-words-to-embeddings/learning.csv};
\addplot[gray,very thick,mark=diamond*] table[x=docs,y=embeddings]{figdata/ml/13-text-from-bag-of-words-to-embeddings/learning.csv};
\addplot[black,thick,mark=o] table[x=docs,y=events]{figdata/ml/13-text-from-bag-of-words-to-embeddings/learning.csv};
\legend{effet injecté,dictionnaire,tf-idf en mots,tf-idf avec bigrammes,mots supervisés,plongements,règles d'événements}
\end{semilogxaxis}
\end{tikzpicture}

\omcaption{Coefficient d'information de rang sur la moitié de test en fonction du nombre de titres étiquetés sur lesquels chaque modèle a été entraîné (le dictionnaire
n'en utilise aucun). Avec 300 titres, aucun mot ne passe le crible supervisé et son score est nul. Données~:
\texttt{ml\_text.learning\_curve}.}\label{fig:ml:text:learning}
\end{omfigure}

Trois enseignements se dégagent des courbes. Un dictionnaire ne coûte aucune étiquette~: le modèle tf-idf en mots ne le dépasse qu'entre
1\,000 et 3\,000 titres étiquetés, la liste de mots supervisée et le modèle à bigrammes seulement au-delà de 3\,000, et pour un texte réel, avec
un vocabulaire bien plus vaste et des signaux plus faibles, disposer d'assez d'étiquettes revient à attendre des années de dépêches. La \omterm{def:ml:linear-and-regularised-baselines:logistic}{régression logistique} sur les mots tf-idf (le modèle du chapitre~4, avec une pénalité ridge) se
place en tête des scores du \omterm{def:ml:text-from-bag-of-words-to-embeddings:bow}{sac de mots}, à 0.273 contre 0.224 pour le dictionnaire. Quant aux bigrammes, qui peuvent représenter les titres composés
qu'un modèle de mots ne voit pas, ils coûtent davantage en variance qu'ils n'apportent~: 0.257 contre 0.273, avec près de quatre fois
plus de variables et des composés qui ne forment que 9~\% des titres. Une représentation plus riche a besoin de plus de données pour se rentabiliser.

\section{Thèmes et événements}

\begin{definition}[Modèle thématique, allocation de Dirichlet latente]\label{def:ml:text-from-bag-of-words-to-embeddings:lda}
Un \emph{modèle thématique}\index{modèle thématique} décrit chaque document comme un mélange de quelques thèmes, chaque thème étant une distribution sur
les mots, le tout estimé à partir des effectifs de mots sans étiquettes. L'\emph{allocation de Dirichlet latente}\index{allocation de Dirichlet latente}
(LDA) est le modèle thématique bayésien dans lequel la part de chaque thème dans un document et les probabilités de mots de chaque thème suivent des
a priori de Dirichlet, et chaque mot est tiré en commençant par tirer son thème selon les parts du document (Blei, Ng et Jordan, 2002).
\end{definition}

\begin{table}[htbp]
\centering\small
\begin{tabular}{@{}cl@{}}
\toprule
thème & les six mots les plus probables\\
\midrule
1 & analysts, fell, short, expectations, forecasts, exceeded\\
2 & outlook, guidance, but, expectations, forecasts, estimates\\
3 & annual, meeting, record, attendance, reports, holds\\
4 & full, forecast, year, reports, line, costs\\
5 & wins, authorises, program, repurchase, industry, award\\
6 & announces, gets, rating, earnings, schedules, call\\
\bottomrule
\end{tabular}
\caption{Six thèmes ajustés par LDA aux titres d'entraînement (le marqueur de société, les prépositions et les mots de contexte
retirés). Données~: \texttt{ml\_text.topics}.}\label{tab:ml:text:topics}
\end{table}

Le \Cref{tab:ml:text:topics} montre ce qu'un \omterm{def:ml:text-from-bag-of-words-to-embeddings:lda}{modèle thématique} trouve dans des titres de six mots~: les gabarits du corpus, et non ses actualités.
Les dépassements et les manquements partagent le thème~1, les changements de guidance et les titres combinant l'un ou l'autre sens partagent le thème~2, et un rachat d'actions partage
le thème~5 avec un prix sectoriel~; en
moyenne, seuls 63~\% des titres d'un événement tombent dans son thème le plus fréquent. Un thème est un motif de cooccurrence, et des mots
de sens opposé cooccurrent avec les mêmes mots. Les thèmes sont utiles pour organiser de longs documents (quelles sections
d'un dépôt ont changé, ce à quoi une conférence de résultats a consacré son temps) et comme variables pour un modèle supervisé, et non comme un indicateur de sentiment.

\begin{definition}[Extraction d'événements]\label{def:ml:text-from-bag-of-words-to-embeddings:events}
L'\emph{extraction d'événements}\index{extraction d'événements} (event extraction) transforme un texte en enregistrements structurés (qui, quel type d'événement, quel
sens, quand), au moyen de règles sur des mots-clés et de la syntaxe ou par un étiqueteur entraîné, afin qu'un modèle puisse être estimé sur des événements plutôt
que sur des mots.
\end{definition}

L'extracteur du chapitre tient en une douzaine de règles~: un mot-clé par événement, «~not~» devant une baisse de guidance, un résultat et un
changement de guidance reliés par «~but~». Il lit correctement chacun des 20\,000 titres, et un score qui affecte à chaque
événement extrait son rendement moyen d'entraînement atteint 0.292, contre 0.295 pour l'effet injecté, au-dessus de tous les modèles de \omterm{def:ml:text-from-bag-of-words-to-embeddings:bow}{sac de mots}.
Il le devrait~: ses règles ont été écrites à partir des gabarits propres au corpus. Sur un titre pour lequel il n'a pas été conçu («~surpasses
consensus~») il ne renvoie rien. Les règles sont précises et fragiles~; les modèles appris sont robustes et avides d'étiquettes~; les systèmes de production
utilisent des règles pour les événements les plus importants et des étiqueteurs appris pour le reste, et mesurent les deux sur des textes
à partir desquels ils n'ont pas été construits.

\section{Plongements de mots}

\begin{definition}[Plongement de mots]\label{def:ml:text-from-bag-of-words-to-embeddings:embedding}
Un \emph{plongement de mots}\index{plongement de mots} (word embedding) associe à chaque mot un vecteur dense de quelques dizaines à quelques milliers de nombres réels,
appris de sorte que des mots employés dans des contextes semblables reçoivent des vecteurs proches~: par factorisation d'une matrice de statistiques de cooccurrence, ou
en entraînant un réseau à prédire un mot à partir de ses voisins (Mikolov et coauteurs, 2013).
\end{definition}

Les plongements du chapitre sont du premier type (le \cref{lst:ml:text:ppmi})~: l'information mutuelle ponctuelle positive des mots
apparaissant dans un même titre, réduite à 16 dimensions par une décomposition en valeurs singulières (livre~4, chapitre~25), lignes
ramenées à une norme unitaire. Elles apprennent exactement ce que promet la définition. Les synonymes qui occupent la même place, «~beat~» et
«~tops~», «~raised~» et «~lifts~», «~guidance~» et «~outlook~», ont des similarités cosinus supérieures à 0.99~; il en va de même des antonymes qui
occupent la même place, «~beat~» et «~missed~» à 0.95 (la \cref{fig:ml:text:pairs}). Le contexte décide de la similarité, et un mot et
son contraire partagent leurs contextes. Une régression ridge sur la moyenne des vecteurs de mots d'un titre a alors peu à
exploiter~: IC 0.094. Les plongements entraînés sur un grand corpus généraliste en contiennent davantage que celui de ce jouet, mais la propriété est la même, et
c'est pourquoi les modèles de sentiment bâtis sur des plongements font l'objet d'un \omterm{def:ml:representation-learning:ssl}{ajustement fin} sur des étiquettes (chapitre~10), et pourquoi les modèles de langage du chapitre~14
sont évalués sur la tâche, et non sur leurs vecteurs.

\begin{omfigure}
\begin{tikzpicture}
\begin{axis}[width=10cm,height=7cm,xbar,bar width=7pt,bar shift=0pt,area legend,xmin=0,xmax=1.05,
  xlabel={similarité cosinus des deux vecteurs de mots},ytick={0,1,2,3,4,5,6,7,8,9,10},yticklabels from table={figdata/ml/13-text-from-bag-of-words-to-embeddings/pairs.csv}{pair},y dir=reverse,
  tick label style={font=\small},label style={font=\small},grid=major,grid style={gray!25},
  legend style={font=\footnotesize,at={(0.5,-0.18)},anchor=north,draw=none,fill=none},legend columns=4,
  table/col sep=comma,enlarge y limits=0.06]
\addplot[fill=omDef,draw=none,restrict expr to domain={\thisrow{i}}{0:2}] table[x=cosine,y=i]{figdata/ml/13-text-from-bag-of-words-to-embeddings/pairs.csv};
\addplot[fill=omProp,draw=none,restrict expr to domain={\thisrow{i}}{3:4}] table[x=cosine,y=i]{figdata/ml/13-text-from-bag-of-words-to-embeddings/pairs.csv};
\addplot[fill=omThm,draw=none,restrict expr to domain={\thisrow{i}}{5:8}] table[x=cosine,y=i]{figdata/ml/13-text-from-bag-of-words-to-embeddings/pairs.csv};
\addplot[fill=gray,draw=none,restrict expr to domain={\thisrow{i}}{9:10}] table[x=cosine,y=i]{figdata/ml/13-text-from-bag-of-words-to-embeddings/pairs.csv};
\legend{même sens,même secteur,opposés,sans lien}
\end{axis}
\end{tikzpicture}

\omcaption{Similarités cosinus des \omterm{def:ml:text-from-bag-of-words-to-embeddings:embedding}{plongements de mots} appris à partir des cooccurrences des titres d'entraînement. Les opposés qui occupent
la même place sont aussi proches que les synonymes. Données~: \texttt{ml\_text.pairs}.}\label{fig:ml:text:pairs}
\end{omfigure}

\section{Liaison d'entités aux titres}

\begin{definition}[Reconnaissance d'entités nommées, liaison d'entités]\label{def:ml:text-from-bag-of-words-to-embeddings:ner}
La \emph{reconnaissance d'entités nommées}\index{reconnaissance d'entités nommées} (named-entity recognition) repère les passages d'un texte qui désignent des entités (sociétés,
personnes, lieux, produits) et leur type. \emph{liaison d'entités}\index{liaison d'entités} (entity linking) associe chaque passage reconnu à un
identifiant dans une base de connaissances~; pour une société de trading, à un identifiant permanent du référentiel des titres (livre~7, chapitre~4) à
la date du document.
\end{definition}

Un modèle qui lit parfaitement les dépêches et opère sur la mauvaise société vaut moins que pas de modèle du tout. Le module de reconnaissance du chapitre est un
dictionnaire d'entités~: chaque nom, nom court et ticker qu'une société a jamais porté, appariés du plus long au plus court, ce qui retrouve chaque mention dans le
corpus, car celui-ci n'emploie aucun autre nom~; les reconnaisseurs réels étiquettent aussi les noms jamais vus, à partir de la majuscule,
de la syntaxe et d'un étiqueteur entraîné. Le module de liaison (le \cref{lst:ml:text:link}) résout alors un ticker par le référentiel des titres (security master) et
un nom par une table d'alias, toutes deux à la date du titre, et départage plusieurs candidats grâce aux mots de secteur figurant dans le
titre, puis par la couverture.

\begin{table}[htbp]
\centering\small
\begin{tabular}{@{}lcccccc@{}}
\toprule
& \multicolumn{3}{c}{alias à la date du titre} & \multicolumn{3}{c}{alias actuels}\\
\cmidrule(lr){2-4}\cmidrule(l){5-7}
mention & correct & lié & précision & correct & lié & précision\\
\midrule
nom complet & 100.0 & 100.0 & 100.0 & 89.3 & 89.3 & 100.0\\
nom court & 98.2 & 100.0 & 98.2 & 86.8 & 88.6 & 98.0\\
ticker & 100.0 & 100.0 & 100.0 & 84.5 & 92.4 & 91.5\\
\midrule
ensemble & 99.5 & 100.0 & 99.5 & 87.5 & 89.7 & 97.5\\
\bottomrule
\end{tabular}
\caption{\omterm{def:ml:text-from-bag-of-words-to-embeddings:ner}{Liaison d'entités} sur les 20\,000 titres (\%)~: part liée à la bonne société, part liée à une société quelconque, et
précision (bonnes parmi les liées). Données~: \texttt{ml\_text.linking}.}\label{tab:ml:text:linking}
\end{table}

Le \Cref{tab:ml:text:linking} donne le second résultat du chapitre. Avec des alias valables à chaque date, le linker a raison dans 99.5~\%
des cas~; ses seules erreurs portent sur le nom court Apex dans les 245 dépêches qui ne comportent aucun mot de secteur, où il se rabat
sur l'Apex le plus couvert et a raison dans 58.0~\% des cas (avec un mot de secteur, il a raison dans les 737). Avec les alias d'aujourd'hui,
il a raison dans 87.5~\% des cas~: les anciens noms et les anciens tickers ne renvoient à rien, et les quatre tickers réutilisés renvoient vers la nouvelle cotation,
de sorte que 8.5~\% des mentions de ticker qu'il relie pointent vers la mauvaise société, sans rien signaler. Un backtest sur ces liens négocierait
une dépêche de rachat concernant une société sur les actions d'une autre~; la règle du livre~7, chapitre~3, s'applique aux alias comme aux
cours.

\begin{method}[Un pipeline de texte pour le trading]\label{met:ml:text:method}
\begin{enumerate}
\item Reconnaître les mentions et les relier, à la date du document, à des identifiants permanents~; mesurer la précision de la liaison sur un
échantillon étiqueté, par type de mention.
\item Masquer les mentions, tokeniser et construire la représentation (d'abord les mots en TF-IDF).
\item Partir d'un dictionnaire comme référence~; ajuster des scores supervisés sur les rendements dans des plis walk-forward, avec la fenêtre de réaction
et l'horizon fixés à l'avance.
\item Extraire les événements qui comptent à l'aide de règles ou d'un étiqueteur, et les tester face au modèle en \omterm{def:ml:text-from-bag-of-words-to-embeddings:bow}{sac de mots}.
\item Traiter les plongements et les thèmes comme des variables à valider, et non comme du sentiment.
\end{enumerate}
\end{method}

\section{Tutoriel~: quel Apex~?}\label{tut:ml:text-from-bag-of-words-to-embeddings:tut}

\begin{tutorial}
\textbf{Objectif.} Générer le corpus, reconnaître et lier les mentions à la date du document, et comparer cinq scores de texte par leur
coefficient d'information. \textbf{État final~:} \cref{tab:ml:text:ic,tab:ml:text:topics,tab:ml:text:linking},
\cref{fig:ml:text:learning,fig:ml:text:pairs}.
\begin{tutsteps}
\item \textbf{Une liste de mots supervisée par les rendements.}
\omcode{code/firm/textml/firm_textml.py}{51}{70}{Un filtre de mots supervisé par les rendements.}\label{lst:ml:text:supervised}
\item \textbf{Plongements à partir des cooccurrences.}
\omcode{code/firm/textml/firm_textml.py}{78}{99}{Plongements de mots à partir de l'information mutuelle ponctuelle positive et d'une décomposition en valeurs singulières.}\label{lst:ml:text:ppmi}
\item \textbf{Liaison à la date du document.}
\omcode{code/firm/textml/firm_textml.py}{211}{222}{Liaison d'une mention à un identifiant permanent à une date donnée.}\label{lst:ml:text:link}
\item \textbf{Lancer} \texttt{ml\_text.scores()}, \texttt{learning\_curve()}, \texttt{topics()}, \texttt{pairs()},
\texttt{linking()} et \texttt{fig\_text.py}.
\end{tutsteps}
\textbf{Pour aller plus loin.} Ajouter des formulations à la période de test que la période d'apprentissage n'a jamais vues, et observer comment les règles et
les modèles supervisés échouent différemment~; donner au linker un modèle de contexte à la place d'une liste de mots de secteur.
\end{tutorial}

\section{Construction~: variables textuelles et liaison d'entités}\label{bld:ml:text-from-bag-of-words-to-embeddings:textml}

\begin{build}
\bfield{Rôle} Du texte transformé en scores qui peuvent être validés comme toute autre variable, et rattachés au bon titre.
\bfield{Interface} \texttt{tokenize}, \texttt{ngrams}, \texttt{dictionary\_score(tokens, pos, neg)},
\texttt{supervised\_words(docs, returns, min\_count, z)}, \texttt{supervised\_score},
\texttt{ppmi\_svd(docs, k, min\_count)}, \texttt{doc\_vectors}, \texttt{extract\_event(tokens)}~; \texttt{AliasTable},
\texttt{recognise(text, gazetteer)}, \texttt{link(surface, date, context, table, sm, sectors, prior)}~; \texttt{news\_corpus(n\_docs, seed, days)}
sur \texttt{firm.secmaster}.
\bfield{Règles} Les mentions sont masquées avant que tout modèle ne voie le texte~; les alias et les tickers sont résolus à la date du
document~; les scores supervisés sont ajustés sur des documents antérieurs à ceux qu'ils évaluent.
\bfield{Tests d'acceptation} \texttt{code/firm/textml/tests/}~: jetons et \omterm{def:ml:text-from-bag-of-words-to-embeddings:bow}{n-grammes} d'une dépêche connue~; le filtre supervisé
retrouve un mot injecté et ignore un mot neutre~; les plongements rapprochent les mots aux contextes identiques~; les règles lisent
les composés et les négations~; la liaison à la date suit un changement de ticker, un renommage et un ticker réutilisé, alors que la table du jour ne
le fait pas.
\bfield{Pour aller plus loin} Un étiqueteur de séquences entraîné pour la reconnaissance~; une liaison par modèle de contexte sur les descriptions des candidats~;
le \omterm{def:ml:text-from-bag-of-words-to-embeddings:lda}{modèle thématique} de Ke--Kelly--Xiu au-dessus du filtre.
\end{build}

\begin{omsources}
\item P.~C.~Tetlock, ``Giving content to investor sentiment: the role of media in the stock market'', \emph{Journal of
Finance} 62(3), 2007.
\item T.~Loughran et B.~McDonald, ``When is a liability not a liability? Textual analysis, dictionaries, and 10-Ks'',
\emph{Journal of Finance} 66(1), 2011.
\item Z.~T.~Ke, B.~T.~Kelly et D.~Xiu, ``Predicting returns with text data'', NBER Working Paper 26186, 2019.
\item M.~Gentzkow, B.~Kelly et M.~Taddy, ``Text as data'', \emph{Journal of Economic Literature} 57(3), 2019.
\item D.~M.~Blei, A.~Y.~Ng et M.~I.~Jordan, ``Latent Dirichlet allocation'', \emph{Advances in Neural Information Processing
Systems} 14, MIT Press, 2002.
\item T.~Mikolov, K.~Chen, G.~Corrado et J.~Dean, ``Efficient estimation of word representations in vector space'',
arXiv:1301.3781, 2013.
\end{omsources}

\section{Exercices}

\begin{exercise}[$\star$]\label{exo:ml:text-from-bag-of-words-to-embeddings:1}
Lister les jetons et les bigrammes de «~\texttt{<co>} did not cut guidance~». Quel bigramme porte la négation~?
\end{exercise}

\begin{exercise}[$\star$]\label{exo:ml:text-from-bag-of-words-to-embeddings:2}
Dans un corpus de 20\,000 documents, «~guidance~» apparaît dans 3\,000 et «~halves~» dans 200. Quelles sont leurs fréquences
inverses de document ($\log N/n_w$)~? Lequel pèse le plus dans un document qui contient les deux une fois~?
\end{exercise}

\begin{exercise}[$\star$]\label{exo:ml:text-from-bag-of-words-to-embeddings:3}
Calculer le score de dictionnaire de «~\texttt{<co>} beat estimates but cut guidance~» (six jetons) et de «~\texttt{<co>} wins
industry award~» (quatre jetons) avec les listes du chapitre. Quels sont les effets injectés~?
\end{exercise}

\begin{exercise}[$\star\star$]\label{exo:ml:text-from-bag-of-words-to-embeddings:4}
Pourquoi «~beat~» et «~missed~» reçoivent-ils presque le même plongement ici, et qu'est-ce qui les distinguerait~?
\end{exercise}

\begin{exercise}[$\star\star$]\label{exo:ml:text-from-bag-of-words-to-embeddings:5}
Le filtre supervisé conserve «~to~» et «~be~». Est-ce un bogue~? Dans quels cas le deviendrait-il~?
\end{exercise}

\begin{exercise}[$\star\star$]\label{exo:ml:text-from-bag-of-words-to-embeddings:6}
\emph{Trouver la faille.} «~Nous avons relié cinq ans d'actualités à des tickers avec le fichier de symboles actuel du fournisseur~; 97~\% des liens que nous
avons vérifiés à la main étaient justes, donc la liaison ne pose pas de problème.~»
\end{exercise}

\begin{exercise}[$\star\star\star$]\label{exo:ml:text-from-bag-of-words-to-embeddings:7}
\emph{Programmation.} Entraîner le filtre supervisé et la \omterm{def:ml:linear-and-regularised-baselines:logistic}{régression logistique} sur TF-IDF sur la moitié d'entraînement avec les
mentions de sociétés laissées dans le texte (non masquées). Qu'apprennent les modèles sur les sociétés, et que devient l'IC de test~? Expliquer.
\end{exercise}

\begin{exercise}[$\star\star\star$]\label{exo:ml:text-from-bag-of-words-to-embeddings:8}
Montrer qu'avec les jetons d'un document comme \omterm{def:ml:text-from-bag-of-words-to-embeddings:bow}{sac de mots}, aucun modèle linéaire sur les comptages de mots ne peut attribuer à «~beat \dots\ but cut~» un
effet de $-1.2$ et à «~missed \dots\ but raised~» un effet de $+1.2$ tout en donnant à «~beat~», «~missed~», «~raised~» et «~cut~»
isolément, les effets $+1$, $-1$, $+1.5$ et $-1.5$. Quelles variables le permettraient~?
\end{exercise}

\section{Problème~: Quel Apex~?}

\begin{problem}[{Problème du week-end --- du texte à la bonne transaction}]\label{pb:ml:text-from-bag-of-words-to-embeddings:1}
Le corpus du chapitre~: 20\,000 titres d'actualité, 32 sociétés (trois s'appellent Apex), six changements de ticker, trois renommages, quatre
tickers réutilisés.

\medskip\noindent\textbf{Partie~I --- Représentation.}
\begin{enumerate}
\item Pourquoi masquer les mentions de sociétés avant de construire les variables~?
\item Que fait le TF-IDF à un mot qui apparaît dans chaque titre~?
\item Combien de variables les bigrammes ajoutent-ils, à peu près, et pourquoi n'ont-ils pas aidé ici~?
\item Que dit l'IC de 0.295 de l'effet injecté, sur le meilleur résultat qu'un modèle de texte puisse atteindre~?
\end{enumerate}

\medskip\noindent\textbf{Partie~II --- Scores.}
\begin{enumerate}[resume]
\item Classer les six scores du \cref{tab:ml:text:ic}.
\item En dessous de combien de titres étiquetés le dictionnaire est-il le meilleur score~?
\item Pourquoi les règles événementielles battent-elles tous les modèles appris, et pourquoi ne faut-il pas y croire~?
\item Pourquoi les plongements obtiennent-ils un score de 0.094~?
\end{enumerate}

\medskip\noindent\textbf{Partie~III --- Liaison.}
\begin{enumerate}[resume]
\item Quelle part des mentions la liaison point-in-time identifie-t-elle correctement, et où sont ses erreurs~?
\item Que fait bien la table du jour, et quels sont ses deux modes de défaillance~?
\item Quelle défaillance est la plus grave pour un backtest, une mention non liée ou une mention mal liée~?
\item Quelle est l'exactitude de l'outil de liaison sur Apex, avec et sans un mot sectoriel~?
\end{enumerate}

\medskip\noindent\textbf{Partie~IV --- Le verdict.}
\begin{enumerate}[resume]
\item Énoncer le \emph{résultat nommé}~: l'exactitude et la précision de la \omterm{def:ml:text-from-bag-of-words-to-embeddings:ner}{liaison d'entités}, avec et sans alias point-in-time,
ainsi que l'IC des modèles supervisés par les rendements, face à celui du dictionnaire.
\item Que faut-il construire en premier pour un desk d'actualités~: un meilleur score de sentiment ou une meilleure liaison~?
\item Comment vérifier les \omterm{def:ml:text-from-bag-of-words-to-embeddings:ner}{liaisons d'entités} fournies par un fournisseur~?
\item Comment garantir l'absence de biais temporel dans un modèle de texte supervisé~?
\item À quoi sert un \omterm{def:ml:text-from-bag-of-words-to-embeddings:lda}{modèle thématique} dans ce métier~?
\item Qu'est-ce qui change avec les documents longs (dépôts réglementaires, transcriptions d'appels) plutôt qu'avec les titres~?
\item Que faudrait-il pour que le dictionnaire redevienne compétitif~?
\item En une phrase~: quelle est l'étape la plus importante entre un titre et une transaction~?
\end{enumerate}
\end{problem}

\section{Questions d'entretien}

\begin{interviewq}[$\star$ \iqroles{chercheur, ingénieur ML}]\label{iq:ml:text-from-bag-of-words-to-embeddings:1}
Qu'est-ce que le TF-IDF, et en quoi aide-t-il un modèle linéaire sur du texte~?
\end{interviewq}

\begin{interviewq}[$\star$ \iqroles{chercheur}]\label{iq:ml:text-from-bag-of-words-to-embeddings:2}
Pourquoi existe-t-il des dictionnaires propres à la finance~?
\end{interviewq}

\begin{interviewq}[$\star\star$ \iqroles{chercheur}]\label{iq:ml:text-from-bag-of-words-to-embeddings:3}
Comment construiriez-vous un score de sentiment supervisé par les rendements, et quelles fuites d'information pourraient s'y glisser~?
\end{interviewq}

\begin{interviewq}[$\star\star$ \iqroles{ingénieur ML}]\label{iq:ml:text-from-bag-of-words-to-embeddings:4}
Pourquoi les antonymes sont-ils proches dans l'espace des \omterm{def:ml:text-from-bag-of-words-to-embeddings:embedding}{plongements de mots}, et qu'est-ce que cela implique pour l'usage des plongements comme variables~?
\end{interviewq}

\begin{interviewq}[$\star\star$ \iqroles{ingénieur ML, développeur}]\label{iq:ml:text-from-bag-of-words-to-embeddings:5}
Concevez l'étape de \omterm{def:ml:text-from-bag-of-words-to-embeddings:ner}{liaison d'entités} pour un flux d'actualités qui pilote le trading. Qu'est-ce qui peut mal tourner~?
\end{interviewq}

\begin{interviewq}[$\star\star\star$ \iqroles{chercheur}]\label{iq:ml:text-from-bag-of-words-to-embeddings:6}
L'IC de votre modèle de texte a été divisé par deux lors de sa première année en production. Énumérez les causes que vous vérifieriez, dans l'ordre.
\end{interviewq}
